\section{Cote 134-8: two extensions using purity}
\label{sec:SM-cote1348-purity}

Retain the precise connective setup of BR-13403 and the ringoid
construction of BR-13405. Thus $k$ is commutative,
$R$ is a small $k$-linear category with projective scalar modules
$R(x,y)$, $\mathcal M=\operatorname{Mod}\!-\!R$, and
\[
 \varphi(X)=D\otimes_R X,\qquad
 \Lambda(X)=B_\bullet\otimes_R X,\qquad
 (GC)(x)=\operatorname{Hom}_{\mathrm{Ch}_{\geq0}(k)}(B_\bullet(x),C).
 \tag{SM24.1}
\]
Here $B_\bullet\to D$ is the chosen resolution by coproducts of
left representables. The canonical unit normalization and the full
interval condition, including degree zero, remain in force.
For precision we also write the scalar compatibility used below:
\[
 J\otimes_{\mathcal M}(X\otimes_k U)
   \simeq (J\otimes_{\mathcal M}X)\otimes_k U,\qquad
 \alpha_{X\otimes_k U}=\alpha_X\otimes_k U.
 \tag{SM24.2}
\]
It holds for a $k$-linear tensor product preserving colimits in each
variable: present $U$ by free modules and use coproducts and cokernels.
We make it explicit rather than infer it from an unspecified tensor
operation. It is the scalar compatibility of the tensor setup being
used.

\paragraph{Finding F-SM-24 (the ordinary injective-dimension hypothesis can be replaced).}
Under these hypotheses, either one of the following two conditions
suffices for the specified adjunction to induce
\[
 \mathcal M[W^{-1}]\simeq D_{\geq0}(k),
 \qquad W=\Lambda^{-1}(\mathrm{qis}).
 \tag{SM24.3}
\]
\begin{enumerate}
\item There is a finite $r$ such that every projective $k$-module
      has a pure-injective resolution of length at most $r$.
\item Every $R(x,y)$ is finitely generated projective over $k$.
\end{enumerate}
The second condition imposes no injective or pure-injective dimension
condition on $k$. The first, by the known cardinality theorem stated
below, applies to every countable scalar ring and more generally to
$|k|=\aleph_t$ for a finite integer $t\geq0$.
In the first condition, $R(x,y)$ may be infinitely generated.

\begin{proof}
\emph{Purity conventions and elementary facts.}
A monomorphism of right modules is pure if tensoring with every
left module keeps it injective. Equivalently, the associated short
exact sequence remains exact after applying $\operatorname{Hom}(F,-)$
for every finitely presented right module $F$.
This equivalence also holds for ringoid modules. Indeed, a presentation
of $F$ by finite sums of representables makes a lift of a map
$F\to C$ precisely a solution of a finite system of module equations
whose constants are in the submodule. Tensoring with the cokernel
of the transposed finite matrix tests solvability of that same system.
Arbitrary tensor tests reduce to finitely presented tests because
every module is a filtered colimit of finitely presented modules.
In particular, products of pure short exact sequences are pure:
maps from a finitely presented module into a product are tuples of
maps, and products of surjections of abelian groups are surjective.

A module $E$ is pure injective if maps into $E$ extend across pure
monomorphisms. Products of pure-injective modules are pure injective,
by extending the coordinate maps. Every scalar module $Q$ has a
pure monomorphism into a pure-injective module. Here is the needed
construction, avoiding any envelope assumption. Write
$Q^+=\operatorname{Hom}_{\mathbb Z}(Q,\mathbb Q/\mathbb Z)$.
For every module $V$, the character module $V^+$ is pure injective:
use
\[
 \operatorname{Hom}_k(M,V^+)
   =\operatorname{Hom}_{\mathbb Z}(V\otimes_k M,\mathbb Q/\mathbb Z)
\]
and injectivity of $\mathbb Q/\mathbb Z$.
The evaluation map $e_Q:Q\to Q^{++}$ is pure.
Indeed $e_Q^+:Q^{+++}\to Q^+$ has the section $e_{Q^+}$.
For every $V$, the character dual of $V\otimes e_Q$ is therefore
surjective. Since $\mathbb Q/\mathbb Z$ is an injective cogenerator
of abelian groups, $V\otimes e_Q$ is injective. Thus $Q^{++}$ is
the required pure-injective target.

We will also use the fact that all functors
$\operatorname{Tor}_n^R(D,-)$ preserve pure short exact sequences.
For ringoid modules, the character dual is taken at every object
and reverses the variance; the Ext groups below are in left
$R$-modules.
For completeness, suppose $0\to X\to Y\to Z\to0$ is pure.
The evaluation pairing
$X^+\otimes_R X\to\mathbb Q/\mathbb Z$ extends across
$X^+\otimes_R X\to X^+\otimes_R Y$, which is injective by purity.
The resulting map $X^+\to Y^+$ splits restriction
$Y^+\to X^+$. Thus the character-dual sequence of left
$R$-modules splits.
The natural identity
\[
 \operatorname{Tor}_n^R(D,X)^+
   \simeq \operatorname{Ext}_R^n(D,X^+)
 \tag{SM24.4}
\]
follows by using the projective resolution $B_\bullet$ and the
exactness of character duality.
Applying $\operatorname{Ext}_R^n(D,-)$ to the split dual sequence
and dualizing back proves the assertion. Exactness is reflected
by the character dual because it is an injective cogenerator.
In particular, an $L\varphi$-acyclic module has only
$L\varphi$-acyclic pure submodules.

\emph{Relative pure injectives are killed by the actual interval.}
For $x\in R$ and a scalar module $E$, put
\[
 C_x(E)(y)=\operatorname{Hom}_k(R(x,y),E).
 \tag{SM24.5}
\]
This is the right adjoint to evaluation at $x$:
\[
 \operatorname{Hom}_{\mathcal M}(X,C_x(E))
     =\operatorname{Hom}_k(X(x),E).
 \tag{SM24.6}
\]
Call a monomorphism objectwise pure if its evaluation at every
$x$ is pure over $k$. When $E$ is pure injective, (SM24.6)
shows that $C_x(E)$ is injective with respect to objectwise pure
monomorphisms. So is every product of such modules.

By (SM24.2), the interval monomorphism
$\alpha_X:X\to J\otimes_{\mathcal M}X$ is objectwise pure:
after tensoring its evaluation with any $U$, it is the evaluation
of the assumed monomorphism $\alpha_{X\otimes_k U}$.
Consequently, if
\[
 I=\prod_{a\in A} C_{x_a}(E_a),\qquad E_a\text{ pure injective},
\]
the identity on $I$ extends across $\alpha_I$.
Thus $I$ is a retract of $J\otimes_{\mathcal M}I$.
The full interval acyclicity gives
\[
 L\varphi(I)=0.
 \tag{SM24.7}
\]
No assertion that $I$ is an ordinary injective object is needed.

\emph{The counit test is a product of coinduced scalar coefficients.}
Let $P$ be an acyclic projective object of
$\mathrm{Ch}_{\geq0}(k)$. By BR-13402,
\[
 P=\bigoplus_{n\geq0}D^{n+1}(Q_n)
   =\prod_{n\geq0}D^{n+1}(Q_n),
 \qquad Q_n\text{ projective over }k.
\]
The equality with the product holds degreewise, since only two
disks contribute in any degree. Since $G$ is a right adjoint,
\[
 GP=\prod_{n\geq0}\operatorname{Hom}_k(B_n(-),Q_n).
 \tag{SM24.8}
\]
Writing $B_n=\bigoplus_{a\in A_n}R(x_a,-)$ gives
\[
 GP=\prod_{n\geq0}\prod_{a\in A_n} C_{x_a}(Q_n).
 \tag{SM24.9}
\]
The chain-map formula behind (SM24.8) is
$\operatorname{Hom}_{\mathrm{Ch}}(B,D^{n+1}Q)
=\operatorname{Hom}_k(B_n,Q)$.
In particular, no preservation of coproducts by $G$ was assumed.

\emph{First sufficient condition.}
Choose pure-injective resolutions of length $r$ for all $Q_n$.
The functor $C_x$ is exact, since $R(x,y)$ is projective over $k$
for each $y$. Apply these functors and take the products in
(SM24.9). Products are exact in this module category, so the
result is a finite exact sequence
\[
 0\longrightarrow GP\longrightarrow I^0\longrightarrow
       I^1\longrightarrow\cdots\longrightarrow I^r
       \longrightarrow0
\]
whose terms have the form appearing in (SM24.7).
Every $L\varphi(I^j)$ vanishes.
The long exact sequences of the derived functors, or successive
distinguished triangles, now give $L\varphi(GP)=0$.
This induction is finite; no infinite totalization is involved.

\emph{Second sufficient condition.}
Embed each $Q_n$ purely into $E_n=Q_n^{++}$.
For any finitely presented right $R$-module $F$, the scalar
module $F(x)$ is finitely presented over $k$: evaluate a finite
presentation of $F$ by representables and use the present
finitely generated projectivity of all $R(x,y)$.
If $0\to Q_n\to E_n\to V_n\to0$ is the resulting pure sequence,
(SM24.6) shows that
\[
 0\to C_x(Q_n)\to C_x(E_n)\to C_x(V_n)\to0
\]
is pure as a sequence of right $R$-modules. Exactness uses
projectivity of $R(x,y)$; the purity test uses finite presentation
of $F(x)$ and scalar purity.
Taking products gives a pure monomorphism
\[
 GP\longrightarrow I=
       \prod_{n\geq0}\prod_{a\in A_n}C_{x_a}(E_n).
\]
By (SM24.7), $L\varphi(I)=0$.
By (SM24.4), every $\operatorname{Tor}_n^R(D,GP)$ injects
into the corresponding zero group for $I$.
Thus $L\varphi(GP)=0$ in this case also.

In either case the exact residual condition BR-13403 is satisfied.
Together with the unchanged canonical unit normalization,
F-SM-17 gives (SM24.3) for the specified connective adjunction.
\end{proof}

\paragraph{The countable-ring consequence uses a known theorem.}
Chach\'olski, Neeman, Pitsch, and Scherer (2018),
Theorem~9.2, p.~926, explicitly state the
Kie{\l}pi\'nski--Simson/Gruson--Jensen theorem:
a ring of cardinality $\aleph_t$, $t\in\mathbb N$, has pure
global dimension at most $t+1$. Their preceding paragraph
states this in terms of finite pure-injective dimensions.
Thus F-SM-24(1) applies with $r=t+1$, and in particular with
$r=1$ for every countably infinite scalar ring.
Finite rings are covered as well: apply that theorem to
$k\times\mathbb Z$ and take the first idempotent component of
a pure-injective resolution. A pure sequence and a pure-injective
module over a product ring are exactly pairs of such sequences
and modules over the factors, as follows from the componentwise
tensor and Hom formulas. Hence the resulting resolution over
$k$ has the same length. This proves the stated countable-ring
consequence without claiming a new proof of the cardinality theorem.

\paragraph{Finding F-SM-25 (why the second proof does not automatically extend to arbitrary projective morphism modules).}
Coinduction need not carry a scalar pure short exact sequence to a
pure sequence of ringoid modules when its morphism modules are
infinitely generated projective. The following is a complete
counterexample to that particular inference. It is not a
counterexample to the interval theorem.

\begin{proof}
Take $k=\mathbb Z$ and the one-object ringoid $R=\mathbb Z[t]$,
which is free, but infinitely generated, over $k$.
Let
$E=\widehat{\mathbb Z}=\prod_p\mathbb Z_p$ and $V=E/\mathbb Z$,
where $\mathbb Z$ embeds diagonally.
The exact sequence
\[
 0\longrightarrow\mathbb Z\longrightarrow E
       \longrightarrow V\longrightarrow0
 \tag{SM25.1}
\]
is pure over $\mathbb Z$. For every $n>0$, the map
$\mathbb Z/n\to E/nE$ is the Chinese-remainder isomorphism;
tensoring with $\mathbb Z$ also preserves the injection.
The structure theorem for finitely presented abelian groups,
followed by filtered colimits, gives all tensor tests for purity.

Choose $a=\sum_{j\geq2}2^{j!}\in\mathbb Z_2$.
It is not an integer: its binary expansion has infinitely many
zeros and ones, whereas an ordinary integer has an eventually
zero or eventually one $2$-adic expansion.
Let $e\in E$ have $2$-adic component $a$ and all other components
zero. For each $n\geq0$, choose
$0\leq r_n<2^n$ with $r_n\equiv a\pmod{2^n}$, taking $r_0=0$.
Then
\[
 q_n=(e-r_n)/2^n\in E,\qquad
 2q_{n+1}-q_n=(r_n-r_{n+1})/2^n\in\mathbb Z.
\]
Division is legitimate in the $2$-adic component by the chosen
congruence and in every other component because $2$ is a unit.
Thus $2^{-n}\mapsto q_n+\mathbb Z$ defines a homomorphism
\[
 h:\mathbb Z[1/2]\longrightarrow V.
\]
It has no lift to $E$. The image of $1$ under any such lift would
belong to $\bigcap_n2^n E$, so its $2$-adic component would be
zero. But being a lift would make its difference from $e$ a
diagonal integer, forcing $a$ to be an integer, a contradiction.

Now $F=R/(2t-1)$ is a finitely presented $R$-module whose
underlying abelian group is $\mathbb Z[1/2]$.
For $C(U)=\operatorname{Hom}_{\mathbb Z}(R,U)$, the coinduction
adjunction identifies
\[
 \operatorname{Hom}_R(F,C(U))
       =\operatorname{Hom}_{\mathbb Z}(\mathbb Z[1/2],U).
\]
The coinduced version of (SM25.1) is short exact because
$R$ is free over $\mathbb Z$. It is not pure over $R$:
the displayed Hom map from $C(E)$ to $C(V)$ is not surjective,
by the constructed $h$. This proves the claim.
\end{proof}

\paragraph{The residual generality is stated exactly.}
F-SM-24 removes the ordinary scalar injective-dimension condition
for all countable scalar rings, for the other rings of finite
pure global dimension just described, and for every scalar ring
when the ringoid is locally finitely generated projective.
The remaining case has neither available finite pure-injective
dimension nor the finite-presentation property used in the second
proof. For that case the exact unresolved assertion is still
\[
 \operatorname{Tor}_j^R
 \left(D,\prod_{n\geq0}\prod_{a\in A_n}C_{x_a}(Q_n)\right)=0
 \quad(j\geq0)
 \tag{SM25.2}
\]
for every sequence of projective scalar modules $Q_n$.
F-SM-25 blocks one proposed automatic purity argument for
(SM25.2); it does not provide data satisfying the original
interval and normalization assumptions and violating (SM25.2).
The ordinary finite-injective-dimension proof, the relative-purity
proofs here, and a hypothetical infinite resolution must therefore
not be conflated.

\subsection*{Source for the cardinality input}
Chach\'olski, W., Neeman, A., Pitsch, W., \& Scherer, J. (2018).
Relative homological algebra via truncations.
\emph{Documenta Mathematica, 23}, 895--937.
\url{https://ems.press/content/serial-article-files/26419?nt=1}.
Exact input: \S9.2 and Theorem~9.2, p.~926.
The source attributes this result to
Kie{\l}pi\'nski--Simson (1975) and Gruson--Jensen (1973);
their original proofs are not claimed to have been reconstructed here.

